\section{Bevezetés}

A 3D nyomtatási folyamatok automatizált hibadetekciójának kutatása az elmúlt évtizedben kiemelkedő
figyelmet kapott, különösen az additív gyártás ipari elterjedésével párhuzamosan. A szakirodalom
két fő megközelítési irányt különböztet meg: a \textbf{képalapú} (gépi látás és mélytanulás) és az
\textbf{akusztikai jelalapú} módszereket. Jelen fejezet áttekinti a terület legfontosabb publikációit,
különös tekintettel az FDM nyomtatókra, a valós idejű detekciós rendszerekre, valamint a
konvolúciós neurális hálózatok alkalmazására, amelyek a jelen projekt elméleti alapjait képezik.

% -------------------------------------------------------------------
\section{Képalapú megközelítések}

    \subsection{Klasszikus képfeldolgozás és referencia-alapú anomáliadedekció}

    A képalapú hibadedekció egyik úttörő megközelítése a referenciakép-összehasonlítás: a
    rendszer az élő kamerakép és egy ideális, hibamentes referencia közötti különbségeket keresi.
    Petsiuk és Pearce~\cite{petsiuk2021hog} nyílt forráskódú, rétegenkénti anomáliadetekciós
    rendszert fejlesztett ki, amelyben a referencia képeket nem valódi nyomtatásból, hanem
    G-kód alapján, Blender fizikai renderelőmotorral szintetikusan állítják elő. A detekció
    irányított gradiens hisztogramok (\textit{Histograms of Oriented Gradients}, HOG) hasonlóságának
    mérésén alapul. A módszer hat reprezentatív hibatípus, köztük a spagetti-jelenség, a
    réteg-eltolódás és az idegen tárgy jelenléte, azonosítására terjed ki, és különlegessége,
    hogy előzetes betanítási adathalmazt nem igényel. Tizenkét hasonlóság- és távolságmérték
    összehasonlítása alapján a Pearson-korreláció, a Spearman-féle rangkorreláció és a koszinusz-hasonlóság
    bizonyult a legmegbízhatóbbnak.

    Egy másik referencia-összehasonlításon alapuló megközelítést mutat be Ben Hammouda és
    társai~\cite{benhammouda2024defect} munkája, amelyben az FDM nyomtatás során keletkező hibákat
    képfeldolgozási technikák, él- és textúraanalízis, szegmentálás, segítségével azonosítják.
    A vizsgált módszerek az STL-fájlokból generált ideális képekkel vetik össze a valódi rétegfelvételeket,
    és különböző geometriai eltérések alapján adnak visszajelzést a nyomtatás minőségéről.
    A kísérletsorozat megerősíti, hogy a kétdimenziós összehasonlítás hatékonyan alkalmazható
    olyan makroszintű hibák felismerésére, mint a vetemedés vagy az alulextrúdálás, ám korlátai
    mutatkoznak finom, kis területre lokalizált defektusok esetén.

    \subsection{Mélytanulás alapú, valós idejű hibadedekció}

    A mélytanulási módszerek, különösen a konvolúciós neurális hálózatok (CNN), új lehetőségeket
    nyitottak a 3D nyomtatási hibák valós idejű, automatikus felismerésében.

    Paraskevoudis és társai~\cite{paraskevoudis2020stringing} egy mély konvolúciós neurális hálózatot
    fejlesztett ki kifejezetten a szálazás (\textit{stringing}) hibatípus felismerésére FFF
    nyomtatásoknál. A modellt élő videóbemenetről szerzett, szálazást ábrázoló képeken tanították be,
    és valós nyomtatási környezetben tesztelték. Az eredmények szerint a módszer nagy sebességgel
    és megbízható pontossággal azonosítja a hibát, és lehetőséget nyújt a folyamat automatikus
    megállítására vagy a nyomtatási paraméterek korrekciójára. A tanulmány hangsúlyozza, hogy
    a kameralapú detekció különösen hatékony azon hibatípusoknál, amelyek vizuálisan, a darab
    és folyamat megfigyelése alapján azonosíthatóak, szemben a forrástüneti hibákkal, amelyek
    érzékelőkkel közvetlenül is detektálhatók.

    Zhang és társai~\cite{zhang2023multiaxis} többtengelyes FDM nyomtatáshoz fejlesztett offline
    minőségellenőrző rendszert, amelyben egy CCD-kamera rögzíti az egyes rétegeket, a CNN pedig
    osztályozza a hibákat. Az interlayer stripping (rétegek közötti szétválás) tipikus defektus
    felismerésére validálták a módszert: az osztályozási pontosság 83,1\%, az érzékenység 85,6\%
    volt, a ROC-görbe alatti terület 0,824-et ért el, ami igazolja a módszer megvalósíthatóságát
    és terjeszthetőségét más hibatípusokra is.

    Átfogó áttekintést nyújt a területről Armin és társai~\cite{armin2025review} összefoglaló tanulmánya,
    amely rendszerezi a képfeldolgozás és gépi látás 3D nyomtatási hibadetekcióban alkalmazott
    módszereit. A szerzők külön tárgyalják a hagyományos, kézi jellemzőkinyerésen alapuló és a
    modern, végponttól végpontig automatizált mélytanulási megközelítéseket, valamint a
    teljes referenciás (\textit{full-reference}) és referencia nélküli (\textit{no-reference})
    módszereket. A review következtetése szerint a valós idejű monitorozó rendszerek fejlesztése
    az anyagveszteség csökkentésének és a fenntartható additív gyártásnak kulcseleme, a jövőbeli
    fejlesztések várhatóan az autonóm visszacsatolási mechanizmusok irányába mutatnak.

    \subsection{Háromdimenziós gépi látás}

    A kétdimenziós képalapú módszerek korlátait, különösen kis méretű, mélységi információt
    igénylő defektusok esetén, 3D gépi látással igyekeznek feloldani.

    Zhao és társai~\cite{zhao2023geometric} pontfelhő-alapú detekciós módszert dolgoztak ki 3D
    nyomtatott felületek hibáinak azonosítására. Az általuk javasolt MBH kulcspontdetektor és INRoPS
    jellemzőleíró kombináció először kivon potenciális hibarégiót, majd pontszomszédsági számítással
    precízen meghatározza a hiba alakját. A módszer robusztusabb és pontosabb a korábbi 2D-alapú
    megközelítéseknél, kiemelten kis méretű defektusok felismerésekor, ahol a kétdimenziós
    képelemzés a zaj által elfedett hibaterületeket tévesen minősíti.

% -------------------------------------------------------------------
\section{Akusztikai alapú megközelítések}

Az akusztikai emisszió (\textit{Acoustic Emission}, AE) analízise egy nem romboló vizsgálati módszer,
amelyet a gyártóipar számos területén alkalmaznak a strukturális épség monitorozásához. Az
additív gyártásban a nyomtatás közbeni anyagkölcsönhatások, repedésképződések és rétegkötési
problémák jellegzetes hangjelet keltenek, amelyek elemzéséből a hiba típusára és helyére
következtetni lehet.

Barriga-Machado és társai~\cite{barriga2025acoustic} PLA alapú, extrudálással előállított minták
károsodásmechanizmusait vizsgálták AE-analízissel kombinálva felügyelet nélküli (\textit{DBSCAN},
\textit{k-means}) és hibrid felügyeletes (\textit{CNN} spektrogram alapú) tanulási módszereket.
A DBSCAN két jól elkülöníthető klasztert azonosított: alacsony frekvenciás (0-260\,kHz) és magas
frekvenciás (270-390\,kHz) csoportot, amelyek a filament törési és a rétegközi szétválási
mechanizmusoknak felelnek meg. Az eredmények igazolják, hogy az adathajtott AE-analízis gyors és
nem romboló eszközt jelent az \textit{in-situ} károsodásdetekciókhoz polimer additív gyártásban.

Waheed és Bernadin~\cite{waheed2023acoustic} valós idejű hangjelalapú hibadetektáló rendszert
fejlesztett FDM nyomtatókhoz CNN-modell alkalmazásával. A rendszer a nyomtatás során keletkező
hangjeleket rögzíti és spektrogrammá alakítja, amelyeket a CNN oszt be hibatípus szerint.
A vizsgált hibakategóriák között szerepelt a fúvókadugulás, a filamentszakadás és a fogaskerék-csúszás.
A kontaktus nélküli, skálázható megközelítés, szemben a hagyományos érzékelőalapú megoldásokkal,
nem igényel közvetlen mechanikai integrációt a nyomtatóba, és az előzetes mérések alapján
megbízható valós idejű hibajelzést képes nyújtani.

Bin Zainal és társai~\cite{binzainal2025lpbf} lézeres porfúziós (\textit{Laser Powder Bed Fusion},
L-PBF) fémes additív gyártási folyamatban alkalmaztak akusztikus monitorozást. A mérőrendszer
mikrofonnal rögzítette az L-PBF folyamat során keletkező akusztikus jeleket, amelyekből hullámlet-transzformáció
(\textit{wavelet transform}) segítségével nyertek ki jellemzőket. Véletlenerdő (\textit{random forest})
és legközelebbi szomszéd (\textit{k-nearest neighbor}) osztályozókkal átlagosan 94\%-os
anyagjellemzési pontosságot értek el, igazolva az akusztikai módszer potenciálját a folyamatparaméterek
valós idejű ellenőrzéséhez.

% -------------------------------------------------------------------
\section{YOLO-alapú tárgydetekció ipari alkalmazásokban}

A YOLO (\textit{You Only Look Once}) objektumdetekciós modellek az iparban is egyre szélesebb
körben terjednek, köszönhetően kiemelkedő sebességüknek és pontosságuknak. Jóllehet az alábbi
tanulmány nem 3D nyomtatásra, hanem acélfelületek hibadetekciójára fókuszál, módszertanilag
szorosan kapcsolódik a jelen projektben alkalmazott megközelítéshez.

Wang és társai~\cite{wang2023yolov5steel} valós idejű acélfelület-hibadedekciós rendszert fejlesztett
többléptékű YOLOv5-alapon. A hálózathoz speciális többléptékű feltáró blokkot és térbeli figyelemfigyelmet
(\textit{spatial attention mechanism}) integráltak a különböző méretű felületi hibák hatékony
felismeréséhez. A NEU-DET adathalmazon elért ~72\%-os mAP érték és a valós idejű sebesség
igazolja, hogy a YOLO-architektúra gyári minőség-ellenőrzési feladatokra, a 3D nyomtatás
hibadetekciójához hasonló körülmények között, megbízhatóan alkalmazható. A többléptékű
megközelítés különösen fontos, mivel az additív gyártás hibái is rendkívül változatos méretekben
és megjelenési formákban fordulnak elő.

% -------------------------------------------------------------------
\section{Összefoglalás és a projekt kontextusa}

A fenti szakirodalom áttekintéséből több fontos következtetés vonható le.

Egyrészt, a \textbf{képalapú, mélytanulási módszerek}, különösen a YOLO-architektúrán alapuló
egymenetes detektorok, a valós idejű ipari hibadetekció legsikeresebb eszközei, amelyek a
referenciaalapú és hagyományos képfeldolgozásnál rugalmasabban generalizálnak, és nem igényelnek
előre meghatározott geometriai modellt.

Másrészt, az \textbf{akusztikai megközelítések} komplementer információt nyújtanak, különösen
belső (a kamerával nem látható) strukturális hibák esetén, azonban valós idejű FDM-alkalmazásokban
még korlátozott az elterjedtségük a szükséges szenzorintegráció és a zajszűrés nehézségei miatt.

Harmadrészt, az áttekintett megoldások többsége \textbf{egyetlen kameraállást} alkalmaz, ami
korlátozza a detektálható hibatípusok körét. Ezt a korlátot a jelen projekt egy
\textbf{többkamerás architektúrával} hivatott feloldani, lehetővé téve, hogy a nyomtató különböző
nézőpontjaiból érkező képek együttes feldolgozásával robusztusabb detekciós rendszer valósuljon meg.

A vonatkozó irodalom alapján a projekt a mélytanulás, konkrétan a YOLOv8, és a valós idejű
kamerastreaming kombinációjára épít, igazodva a terület tudományos konszenzusához, miközben az
általánosabb ipari alkalmazhatóságot és skálázhatóságot is szem előtt tartja.
